% =============================================================================
% CAPÍTULO 6 — DISEÑO TÉCNICO Y DETALLADO DE LA BASE DE DATOS
% =============================================================================
\chapter{Diseño Técnico y Detallado de la Base de Datos}
\label{ch:basedatos}
\resumenCapitulo{Este capítulo expone en su totalidad el motor de datos PostgreSQL~15.10
autoalojado en el servidor TIS (administrado vía phpPgAdmin): los esquemas lógicos y el diagrama entidad-relación, el diccionario de datos
exhaustivo por columna, la integridad referencial, la estrategia de indexación, la lógica
programable (\textit{triggers}, funciones y RPC), la concurrencia y las políticas RLS.}

Este capítulo expone en su totalidad el diseño de la base de datos de EthosHub,
autoalojada en el \textbf{servidor TIS de la UMSS} sobre \textbf{PostgreSQL 15.10} y
administrada mediante el cliente web \textbf{phpPgAdmin}
(\url{http://bytebusters.tis.cs.umss.edu.bo/phppgadmin}). Constituye el documento de
referencia para el administrador de bases de datos (DBA): se documenta la arquitectura
de almacenamiento, el modelo entidad-relación, el diccionario de datos exhaustivo
columna por columna, la integridad referencial, la estrategia de indexación, la lógica
programable (\textit{triggers} y funciones PL/pgSQL), el control de concurrencia y la
seguridad del motor mediante políticas de seguridad a nivel de fila (RLS).

\begin{NotaTecnica}
Toda la información de este capítulo fue \textbf{extraída directamente del catálogo de
sistema} de la instancia de producción (\texttt{pg\_catalog}, \texttt{information\_schema})
mediante consultas de introspección. Refleja, por tanto, el estado real del esquema en
la versión 1.1.0 de la plataforma, no un modelo idealizado.
\end{NotaTecnica}

% =============================================================================
\section{Arquitectura del Almacenamiento Relacional}
\label{sec:arq_almacenamiento}

\subsection{Configuración de esquemas lógicos en PostgreSQL}
\label{subsec:esquemas}

La base de datos no concentra todo en el esquema \texttt{public} (de hecho, este se
mantiene vacío de tablas de negocio). En su lugar, EthosHub aplica una
\textbf{segregación por esquemas lógicos} que separa el dominio de aplicación, la
auditoría histórica, las operaciones administrativas y los esquemas gestionados por la
plataforma Supabase. La tabla~\ref{tab:esquemas} resume esta organización.

\vspace{0.2cm}
\noindent
\renewcommand{\arraystretch}{1.3}
\fontsize{8.5}{10.2}\selectfont\sffamily
\begin{tabularx}{\textwidth}{|>{\ttfamily\bfseries\color{colorPrimario}}l|c|X|}
\hline
\rowcolor{headerTabla}
\sffamily\textbf{Esquema} & \sffamily\textbf{Tablas} & \sffamily\textbf{Propósito} \\
\hline
core & 25 & \textbf{Dominio de aplicación}: perfiles, contenido, interacciones, seguridad. Es el esquema principal sobre el que opera el backend con jOOQ. \\
\hline
\rowcolor{grisTecnico}
history & 2 & \textbf{Auditoría histórica} (SCD Tipo 2): conserva el historial completo de cambios de los perfiles. \\
\hline
admin & 1 & \textbf{Operación administrativa}: tabla \texttt{system\_logs} y funciones de métricas, moderación y \textit{logging}. \\
\hline
\rowcolor{grisTecnico}
auth & 23 & Gestionado por \textbf{Supabase Auth}: usuarios, identidades, sesiones, \textit{tokens}. El backend no lo modifica, solo lo referencia. \\
\hline
realtime & 9 & Gestionado por \textbf{Supabase Realtime}: suscripciones y mensajería \textit{push}. \\
\hline
\rowcolor{grisTecnico}
storage & 8 & Gestionado por \textbf{Supabase Storage}: \textit{buckets} y objetos (avatares, PDFs). \\
\hline
vault & 1 (+1 vista) & \textbf{Supabase Vault}: almacenamiento cifrado de secretos. \\
\hline
\end{tabularx}
\label{tab:esquemas}

\begin{AvisoNota}
La separación \texttt{core} / \texttt{history} / \texttt{admin} sigue el principio de
\textbf{mínimo privilegio por dominio}: el código de aplicación opera contra
\texttt{core}; la auditoría se aísla en \texttt{history} (escrita solo por
\textit{triggers}); y las operaciones privilegiadas viven en \texttt{admin}, accesibles
únicamente mediante funciones \texttt{SECURITY DEFINER}.
\end{AvisoNota}

\paragraph{Extensiones habilitadas.} El motor se apoya en un conjunto acotado de
extensiones PostgreSQL, resumidas en la tabla~\ref{tab:extensiones}.

\clearpage
\vspace{0.2cm}
\noindent
\renewcommand{\arraystretch}{1.3}
\fontsize{9}{10.8}\selectfont\sffamily
\begin{tabularx}{\textwidth}{|>{\ttfamily\color{colorPrimario}}l|X|}
\hline
\rowcolor{headerTabla}
\sffamily\textbf{Extensión} & \sffamily\textbf{Uso en EthosHub} \\
\hline
pgcrypto & Funciones criptográficas; generación de UUID con \texttt{gen\_random\_uuid()}. \\
\hline
\rowcolor{grisTecnico}
uuid-ossp & Generación de identificadores universales únicos (UUID). \\
\hline
pg\_trgm & Búsqueda por similitud de trigramas (búsqueda de talento y \textit{skills}). \\
\hline
\rowcolor{grisTecnico}
btree\_gin & Soporte de índices GIN para tipos comunes combinados. \\
\hline
supabase\_vault & Almacenamiento seguro de secretos. \\
\hline
\rowcolor{grisTecnico}
pg\_stat\_statements & Estadísticas de planificación y ejecución de consultas. \\
\hline
\end{tabularx}
\label{tab:extensiones}

\subsection{Diagrama Entidad-Relación (DER) a nivel físico}
\label{subsec:der}

La figura~\ref{fig:der} presenta el DER físico del esquema \texttt{core}, organizado por
agrupaciones de dominio. La entidad raíz es \comando{profiles\_basic} (perfil
profesional), de la que cuelga todo el contenido. La autenticación reside en
\comando{auth.users} (Supabase), referenciada por las dos tablas de perfil.

\vspace{0.2cm}
\begin{figure}[h!]
\centering
\begin{tikzpicture}[
    ent/.style={rectangle, draw=c4Sistema!70!black, fill=c4Componente!25, font=\sffamily\scriptsize\bfseries,
        align=center, minimum width=2.3cm, minimum height=0.65cm, inner sep=2pt},
    raiz/.style={ent, fill=c4Sistema, text=white},
    auth/.style={ent, fill=c4Externo!30, draw=c4Externo!70!black},
    rel/.style={-{Stealth[length=1.6mm]}, semithick, draw=grisISO!80},
    relid/.style={font=\sffamily\tiny, text=grisISO, fill=white, inner sep=1pt}]

    % Núcleo de identidad
    \node[auth] (users) at (0,4.2) {auth.users};
    \node[raiz] (pb) at (-2.4,2.8) {profiles\_basic};
    \node[raiz] (pc) at (2.4,2.8) {profiles\_company};

    % Contenido profesional (izquierda)
    \node[ent] (proj) at (-5.4,1.4) {projects};
    \node[ent] (pmedia) at (-5.4,0.4) {project\_media};
    \node[ent] (work) at (-5.4,-0.6) {work\_experiences};
    \node[ent] (acad) at (-5.4,-1.6) {academic\_records};
    \node[ent] (hskill) at (-3.0,-1.6) {hard\_skills};
    \node[ent] (sskill) at (-3.0,-0.6) {soft\_skills};
    \node[ent] (tags) at (-3.0,-2.6) {global\_skill\_tags};

    % Portafolio
    \node[ent] (pset) at (-0.4,1.4) {portfolio\_settings};
    \node[ent] (pitem) at (-0.4,0.4) {portfolio\_items};
    \node[ent] (cv) at (-0.4,-0.6) {cv\_documents};
    \node[ent] (conn) at (-0.4,-1.6) {profile\_connections};
    \node[ent] (pref) at (-0.4,-2.6) {profile\_preferences};

    % Interacción reclutador (derecha)
    \node[ent] (likes) at (3.2,1.4) {profile\_likes};
    \node[ent] (chat) at (3.2,0.4) {chats};
    \node[ent] (msg) at (5.6,0.4) {chat\_messages};
    \node[ent] (notes) at (3.2,-0.6) {recruiter\_candidate\_notes};
    \node[ent] (views) at (3.2,-1.6) {recruiter\_talent\_views};
    \node[ent] (notif) at (3.2,-2.6) {notifications};

    % Relaciones identidad
    \draw[rel] (pb) -- node[relid]{1:1} (users);
    \draw[rel] (pc) -- node[relid]{1:1} (users);

    % Contenido → profiles_basic
    \foreach \n in {proj,work,acad,hskill,sskill,cv,conn,pref} \draw[rel] (\n) -- (pb);
    \draw[rel] (pmedia) -- node[relid]{N:1} (proj);
    \draw[rel] (hskill) -- node[relid]{N:1} (tags);
    \draw[rel] (pset) -- node[relid]{1:1} (pb);
    \draw[rel] (pitem) -- node[relid]{N:1} (pset);
    \draw[rel] (notif) -- (pb);

    % Interacción → company + basic
    \draw[rel] (likes) -- (pc);
    \draw[rel] (chat) -- (pc);
    \draw[rel] (msg) -- node[relid]{N:1} (chat);
    \draw[rel] (notes) -- (pc);
    \draw[rel] (views) -- (pc);
\end{tikzpicture}
\caption{Diagrama Entidad-Relación físico del esquema \texttt{core} (agrupado por dominio).}
\label{fig:der}
\end{figure}

\begin{NotaTecnica}
\textbf{Patrón de doble perfil.} EthosHub modela dos tipos de usuario en tablas
separadas: \comando{profiles\_basic} (profesionales) y \comando{profiles\_company}
(reclutadores), ambas con PK \texttt{id\_auth} que referencia a \texttt{auth.users(id)}.
Un \textit{trigger} de exclusividad (\texttt{check\_profile\_exclusivity}) garantiza que
un mismo usuario no exista simultáneamente como ambos tipos.
\end{NotaTecnica}

% =============================================================================
% 6.2 DICCIONARIO DE DATOS EXHAUSTIVO
% =============================================================================
\section{Diccionario de Datos Exhaustivo}
\label{sec:diccionario}

Esta sección documenta \textbf{columna por columna} las 25 tablas del esquema
\texttt{core}, agrupadas por dominio funcional. Para cada columna se indica su nombre
técnico, el tipo de dato nativo de PostgreSQL, su nulidad, el valor por defecto y una
descripción técnica de su uso. Las convenciones transversales del esquema son:

\begin{itemize}[noitemsep]
    \item \textbf{Claves primarias UUID:} la mayoría de las tablas usan UUID generados
    con \comando{gen\_random\_uuid()}, evitando claves secuenciales predecibles.
    \item \textbf{Marcas temporales:} \comando{created\_at} y \comando{updated\_at}
    (tipo \texttt{timestamptz}, defecto \texttt{now()}) en casi todas las tablas.
    \item \textbf{Borrado lógico:} la columna \comando{deleted\_at} (\textit{nullable})
    implementa \textit{soft-delete}; un valor no nulo marca el registro como eliminado
    sin borrarlo físicamente.
\end{itemize}

% -----------------------------------------------------------------------------
\subsection{Esquema de Identidad y Cuentas}
\label{subsec:dic_identidad}

El núcleo de identidad lo forman las dos tablas de perfil y sus dependencias de
seguridad, todas ancladas a \texttt{auth.users(id)} de Supabase.

\paragraph{Tabla \texttt{profiles\_basic} (perfil profesional).}
\begin{DiccTabla}{profiles\_basic}
id\_auth & uuid & NO & — & PK. FK a \texttt{auth.users(id)}; identidad del usuario. \\ \hline
email & text & NO & — & Correo electrónico (único entre activos). \\ \hline
first\_name & text & NO & \texttt{''} & Nombre. \\ \hline
last\_name & text & NO & \texttt{''} & Apellido. \\ \hline
phone\_code & varchar(10) & SÍ & — & Prefijo telefónico (p.~ej. \texttt{+591}). \\ \hline
phone\_number & varchar(30) & SÍ & — & Número telefónico sin prefijo. \\ \hline
country\_code & varchar(10) & SÍ & — & Código de país. \\ \hline
bio & text & SÍ & — & Biografía / narrativa profesional. \\ \hline
location & text & SÍ & — & Ubicación textual. \\ \hline
website & text & SÍ & — & Sitio web personal. \\ \hline
avatar\_url & text & SÍ & — & URL de la foto de perfil. \\ \hline
provider & text & NO & \texttt{'email'} & Proveedor de alta (email / OAuth). \\ \hline
is\_active & boolean & NO & \texttt{true} & Cuenta activa. \\ \hline
deleted\_at & timestamptz & SÍ & — & Marca de borrado lógico. \\ \hline
created\_at & timestamptz & NO & \texttt{now()} & Fecha de creación. \\ \hline
updated\_at & timestamptz & NO & \texttt{now()} & Fecha de última modificación. \\ \hline
seniority & varchar(50) & SÍ & — & Nivel de seniority. \\ \hline
availability\_status & varchar(50) & SÍ & — & Estado: Disponible / Ocupado / Incógnito (CHECK). \\ \hline
last\_password\_changed\_at & timestamptz & SÍ & — & Auditoría de cambio de contraseña. \\ \hline
last\_email\_changed\_at & timestamptz & SÍ & — & Auditoría de cambio de correo. \\ \hline
latitude & double precision & SÍ & — & Latitud geográfica. \\ \hline
longitude & double precision & SÍ & — & Longitud geográfica. \\ \hline
oauth\_providers\_linked & text[] & NO & \texttt{'\{\}'} & Proveedores OAuth vinculados. \\
\end{DiccTabla}

\paragraph{Tabla \texttt{profiles\_company} (perfil reclutador).}
\begin{DiccTabla}{profiles\_company}
id\_auth & uuid & NO & — & PK. FK a \texttt{auth.users(id)}. \\ \hline
email & text & NO & — & Correo del reclutador. \\ \hline
first\_name & text & NO & \texttt{''} & Nombre del contacto. \\ \hline
last\_name & text & NO & \texttt{''} & Apellido del contacto. \\ \hline
phone\_code & varchar(10) & SÍ & — & Prefijo telefónico. \\ \hline
phone\_number & varchar(30) & SÍ & — & Número telefónico. \\ \hline
country\_code & varchar(10) & SÍ & — & Código de país. \\ \hline
company\_name & text & SÍ & — & Razón social de la empresa. \\ \hline
company\_website & text & SÍ & — & Sitio web corporativo. \\ \hline
company\_size & text & SÍ & — & Tamaño de la empresa. \\ \hline
company\_description & text & SÍ & — & Descripción de la empresa. \\ \hline
industry & text & SÍ & — & Industria / sector. \\ \hline
avatar\_url & text & SÍ & — & Logo / avatar. \\ \hline
provider & text & NO & \texttt{'email'} & Proveedor de alta. \\ \hline
is\_active & boolean & NO & \texttt{true} & Cuenta activa. \\ \hline
deleted\_at & timestamptz & SÍ & — & Borrado lógico. \\ \hline
created\_at & timestamptz & NO & \texttt{now()} & Creación. \\ \hline
updated\_at & timestamptz & NO & \texttt{now()} & Modificación. \\ \hline
oauth\_providers\_linked & text[] & NO & \texttt{'\{\}'} & Proveedores OAuth vinculados. \\
\end{DiccTabla}

\paragraph{Tabla \texttt{profile\_preferences} (preferencias del usuario).}
\begin{DiccTabla}{profile\_preferences}
profile\_id & uuid & NO & — & PK. FK a \texttt{auth.users(id)}. \\ \hline
language & varchar(5) & NO & \texttt{'es'} & Idioma de la interfaz (es/en/pt). \\ \hline
theme & varchar(10) & NO & \texttt{'dark'} & Tema visual. \\ \hline
show\_github\_heatmap & boolean & NO & \texttt{true} & Mostrar mapa de calor de GitHub. \\ \hline
show\_linkedin\_recommendations & boolean & NO & \texttt{false} & Mostrar recomendaciones de LinkedIn. \\ \hline
section\_order & text[] & NO & (bio, skills, …) & Orden de secciones del portafolio. \\ \hline
notifications & jsonb & NO & \texttt{'\{\}'} & Preferencias de notificación. \\ \hline
privacy & jsonb & NO & \texttt{'\{\}'} & Preferencias de privacidad. \\ \hline
created\_at & timestamptz & NO & \texttt{now()} & Creación. \\ \hline
updated\_at & timestamptz & NO & \texttt{now()} & Modificación. \\
\end{DiccTabla}

% -----------------------------------------------------------------------------
\subsection{Esquema de Contenido Profesional}
\label{subsec:dic_contenido}

\paragraph{Tabla \texttt{projects} (proyectos).}
\begin{DiccTabla}{projects}
project\_id & uuid & NO & \texttt{gen\_random\_uuid()} & PK. \\ \hline
profile\_id & uuid & NO & — & FK a \texttt{profiles\_basic} (CASCADE). \\ \hline
title & varchar(100) & NO & — & Título (no vacío, CHECK). \\ \hline
description & varchar(500) & NO & — & Descripción (no vacía, CHECK). \\ \hline
category & varchar(20) & NO & \texttt{'Other'} & Web/Mobile/API/Data /DevOps/Other (CHECK). \\ \hline
status & varchar(20) & NO & \texttt{'draft'} & draft/in\_progress /completed/archived (CHECK). \\ \hline
visibility & varchar(10) & NO & \texttt{'Private'} & Public/Private (CHECK). \\ \hline
is\_featured & boolean & NO & \texttt{false} & Proyecto destacado. \\ \hline
thumbnail & text & SÍ & — & Miniatura. \\ \hline
repository\_url & text & SÍ & — & URL del repositorio. \\ \hline
role & varchar(100) & NO & \texttt{''} & Rol desempeñado. \\ \hline
technologies & text[] & NO & \texttt{'\{\}'} & Tecnologías empleadas. \\ \hline
start\_date & date & SÍ & — & Fecha de inicio. \\ \hline
end\_date & date & SÍ & — & Fecha de fin (>= inicio, CHECK). \\ \hline
results & text & SÍ & — & Resultados / logros. \\ \hline
created\_at & timestamptz & NO & \texttt{now()} & Creación. \\ \hline
updated\_at & timestamptz & NO & \texttt{now()} & Modificación. \\ \hline
deleted\_at & timestamptz & SÍ & — & Borrado lógico. \\
\end{DiccTabla}

\paragraph{Tabla \texttt{project\_media} (multimedia de proyectos).}
\begin{DiccTabla}{project\_media}
media\_id & uuid & NO & \texttt{gen\_random\_uuid()} & PK. \\ \hline
project\_id & uuid & NO & — & FK a \texttt{projects} (CASCADE). \\ \hline
url & text & NO & — & URL del recurso (no vacía, CHECK). \\ \hline
type & varchar(20) & NO & — & youtube/vimeo/ figma/slides/pdf/document/image (CHECK). \\ \hline
title & varchar(200) & NO & \texttt{''} & Título del recurso. \\ \hline
size & bigint & SÍ & — & Tamaño en bytes (> 0, CHECK). \\ \hline
created\_at & timestamptz & NO & \texttt{now()} & Creación. \\
\end{DiccTabla}

\paragraph{Tabla \texttt{work\_experiences} (experiencia laboral).}
\begin{DiccTabla}{work\_experiences}
work\_experience\_id & uuid & NO & \texttt{gen\_random\_uuid()} & PK. \\ \hline
profile\_id & uuid & NO & — & FK a \texttt{profiles\_basic} (CASCADE). \\ \hline
company\_name & text & NO & — & Empresa (no vacía, CHECK). \\ \hline
job\_title & text & NO & — & Cargo (no vacío, CHECK). \\ \hline
location & text & SÍ & — & Ubicación. \\ \hline
description & text & SÍ & — & Descripción (<= 1000 car., CHECK). \\ \hline
is\_current & boolean & NO & \texttt{false} & Trabajo actual. \\ \hline
is\_freelance & boolean & NO & \texttt{false} & Trabajo independiente. \\ \hline
logo\_url & text & SÍ & — & Logo de la empresa. \\ \hline
company\_image\_url & text & SÍ & — & Imagen de la empresa. \\ \hline
company\_url & text & SÍ & — & Sitio web de la empresa. \\ \hline
start\_date & date & SÍ & — & Inicio. \\ \hline
end\_date & date & SÍ & — & Fin (>= inicio, CHECK). \\ \hline
technologies & text[] & NO & \texttt{'\{\}'} & Tecnologías. \\ \hline
created\_at & timestamptz & NO & \texttt{now()} & Creación. \\ \hline
updated\_at & timestamptz & NO & \texttt{now()} & Modificación. \\ \hline
deleted\_at & timestamptz & SÍ & — & Borrado lógico. \\ \hline
latitude & double precision & SÍ & — & Latitud (geolocalización). \\ \hline
longitude & double precision & SÍ & — & Longitud (geolocalización). \\
\end{DiccTabla}

\paragraph{Tabla \texttt{academic\_records} (formación académica).}
\begin{DiccTabla}{academic\_records}
academic\_record\_id & uuid & NO & \texttt{gen\_random\_uuid()} & PK. \\ \hline
profile\_id & uuid & NO & — & FK a \texttt{profiles\_basic} (CASCADE). \\ \hline
institution & text & NO & — & Institución (no vacía, CHECK). \\ \hline
degree & text & NO & — & Título obtenido (no vacío, CHECK). \\ \hline
field\_of\_study & text & SÍ & — & Campo de estudio. \\ \hline
education\_type & text & NO & \texttt{'university'} & Enum (CHECK): \pathbreak{university/master/phd/certification/course/bootcamp/high_school}. \\ \hline
start\_date & date & NO & — & Inicio. \\ \hline
end\_date & date & SÍ & — & Fin (>= inicio, CHECK). \\ \hline
gpa & numeric(5,2) & SÍ & — & Promedio (0–100, CHECK). \\ \hline
credential\_url & text & SÍ & — & URL de la credencial. \\ \hline
verification\_url & text & SÍ & — & URL de verificación. \\ \hline
institution\_logo\_url & text & SÍ & — & Logo de la institución. \\ \hline
is\_visible & boolean & NO & \texttt{true} & Visibilidad pública. \\ \hline
description & text & SÍ & — & Descripción. \\ \hline
created\_at & timestamptz & NO & \texttt{now()} & Creación. \\ \hline
updated\_at & timestamptz & NO & \texttt{now()} & Modificación. \\ \hline
deleted\_at & timestamptz & SÍ & — & Borrado lógico. \\
\end{DiccTabla}

\paragraph{Tabla \texttt{global\_skill\_tags} (catálogo global de habilidades).}
\begin{DiccTabla}{global\_skill\_tags}
id & uuid & NO & \texttt{gen\_random\_uuid()} & PK. \\ \hline
name & varchar(100) & NO & — & Nombre de la habilidad (único, CHECK). \\ \hline
category & varchar(50) & NO & — & Frontend/Backend/Data/ Infrastructure/Mobile/Design/Soft Skill (CHECK). \\ \hline
is\_normalized & boolean & NO & \texttt{false} & Tag normalizado por moderación. \\ \hline
created\_at & timestamptz & NO & \texttt{now()} & Creación. \\ \hline
updated\_at & timestamptz & NO & \texttt{now()} & Modificación. \\
\end{DiccTabla}

\paragraph{Tabla \texttt{hard\_skills} (habilidades duras del perfil).}
\begin{DiccTabla}{hard\_skills}
id & uuid & NO & \texttt{gen\_random\_uuid()} & PK. \\ \hline
profile\_id & uuid & NO & — & FK a \texttt{profiles\_basic} (CASCADE). \\ \hline
tag\_id & uuid & NO & — & FK a \texttt{global\_skill\_tags} (CASCADE). \\ \hline
level & varchar(20) & NO & — & Junior/Mid/Senior (CHECK). \\ \hline
is\_top & boolean & NO & \texttt{false} & Habilidad destacada. \\ \hline
top\_order & integer & SÍ & — & Orden 1–3 (solo si is\_top, CHECK). \\ \hline
created\_at & timestamptz & NO & \texttt{now()} & Creación. \\ \hline
updated\_at & timestamptz & NO & \texttt{now()} & Modificación. \\ \hline
deleted\_at & timestamptz & SÍ & — & Borrado lógico. \\
\end{DiccTabla}

\paragraph{Tabla \texttt{soft\_skills} (habilidades blandas).}
\begin{DiccTabla}{soft\_skills}
id & uuid & NO & \texttt{gen\_random\_uuid()} & PK. \\ \hline
profile\_id & uuid & NO & — & FK a \texttt{profiles\_basic} (CASCADE). \\ \hline
title & varchar(200) & NO & — & Título (no vacío, CHECK). \\ \hline
description & varchar(250) & SÍ & — & Descripción. \\ \hline
created\_at & timestamptz & NO & \texttt{now()} & Creación. \\ \hline
updated\_at & timestamptz & NO & \texttt{now()} & Modificación. \\ \hline
deleted\_at & timestamptz & SÍ & — & Borrado lógico. \\
\end{DiccTabla}

\paragraph{Tabla \texttt{portfolio\_settings} (configuración del portafolio).}
\begin{DiccTabla}{portfolio\_settings}
portfolio\_settings\_id & uuid & NO & \texttt{gen\_random\_uuid()} & PK. \\ \hline
profile\_id & uuid & NO & — & FK; único por perfil. \\ \hline
is\_published & boolean & NO & \texttt{false} & Portafolio publicado. \\ \hline
slug & text & SÍ & — & Slug público único (índice parcial). \\ \hline
custom\_headline & text & SÍ & — & Titular personalizado. \\ \hline
custom\_bio & text & SÍ & — & Biografía personalizada. \\ \hline
accent\_color & text & NO & \texttt{'\#8b5cf6'} & Color de acento. \\ \hline
show\_email & boolean & NO & \texttt{false} & Mostrar correo. \\ \hline
show\_location & boolean & NO & \texttt{true} & Mostrar ubicación. \\ \hline
show\_website & boolean & NO & \texttt{true} & Mostrar sitio web. \\ \hline
views\_count & integer & NO & \texttt{0} & Contador de visitas. \\ \hline
created\_at & timestamptz & NO & \texttt{now()} & Creación. \\ \hline
updated\_at & timestamptz & NO & \texttt{now()} & Modificación. \\ \hline
seo\_title & varchar(60) & SÍ & — & Título SEO. \\ \hline
seo\_description & varchar(160) & SÍ & — & Descripción SEO. \\ \hline
show\_connections & boolean & NO & \texttt{true} & Mostrar conexiones. \\ \hline
show\_github\_heatmap & boolean & NO & \texttt{false} & Mostrar heatmap GitHub. \\ \hline
selected\_cv\_document\_id & uuid & SÍ & — & FK a \texttt{cv\_documents} (SET NULL). \\ \hline
cv\_pdf\_url & text & SÍ & — & URL del CV en PDF. \\
\end{DiccTabla}

\paragraph{Tabla \texttt{portfolio\_items} (ítems del portafolio).}
\begin{DiccTabla}{portfolio\_items}
portfolio\_item\_id & uuid & NO & \texttt{gen\_random\_uuid()} & PK. \\ \hline
portfolio\_settings\_id & uuid & NO & — & FK a \texttt{portfolio\_settings} (CASCADE). \\ \hline
item\_type & text & NO & — & project/experience/ education/hard\_skill/soft\_skill (CHECK). \\ \hline
item\_id & uuid & NO & — & ID polimórfico del ítem referenciado. \\ \hline
display\_order & integer & NO & \texttt{0} & Orden de presentación. \\ \hline
created\_at & timestamptz & NO & \texttt{now()} & Creación. \\
\end{DiccTabla}

\paragraph{Tabla \texttt{cv\_documents} (documentos de CV).}
\begin{DiccTabla}{cv\_documents}
id & uuid & NO & \texttt{gen\_random\_uuid()} & PK. \\ \hline
profile\_id & uuid & NO & — & FK a \texttt{profiles\_basic} (CASCADE). \\ \hline
title & text & NO & — & Título (1–200 car., CHECK). \\ \hline
mode & text & NO & \texttt{'markdown'} & markdown/latex (CHECK). \\ \hline
content & text & NO & \texttt{''} & Contenido del CV. \\ \hline
template\_id & text & SÍ & — & Plantilla aplicada. \\ \hline
created\_at & timestamptz & NO & \texttt{now()} & Creación. \\ \hline
updated\_at & timestamptz & NO & \texttt{now()} & Modificación (trigger). \\ \hline
deleted\_at & timestamptz & SÍ & — & Borrado lógico. \\
\end{DiccTabla}

\paragraph{Tabla \texttt{profile\_connections} (conexiones externas).}
\begin{DiccTabla}{profile\_connections}
id & uuid & NO & \texttt{gen\_random\_uuid()} & PK. \\ \hline
profile\_id & uuid & NO & — & FK a \texttt{auth.users(id)} (CASCADE). \\ \hline
provider & text & NO & — & +20 proveedores o \texttt{custom} (CHECK). \\ \hline
status & text & NO & \texttt{'connected'} & connected/disconnected /pending (CHECK). \\ \hline
profile\_handle & text & SÍ & — & Handle del usuario en el proveedor. \\ \hline
provider\_url & text & SÍ & — & URL del perfil externo (enfoque URL-first). \\ \hline
access\_token & text & SÍ & — & Token de acceso (si aplica OAuth). \\ \hline
token\_expires\_at & timestamptz & SÍ & — & Expiración del token. \\ \hline
scopes & text & SÍ & — & Permisos concedidos. \\ \hline
last\_synced\_at & timestamptz & SÍ & — & Última sincronización. \\ \hline
api\_health & text & NO & \texttt{'healthy'} & healthy/degraded/down (CHECK). \\ \hline
provider\_metadata & jsonb & SÍ & \texttt{'\{\}'} & Metadatos del proveedor. \\ \hline
created\_at & timestamptz & NO & \texttt{now()} & Creación. \\ \hline
updated\_at & timestamptz & NO & \texttt{now()} & Modificación (trigger). \\ \hline
provider\_kind & text & NO & \texttt{'manual'} & auth/manual/custom (CHECK). \\ \hline
display\_label & text & SÍ & — & Etiqueta a mostrar. \\
\end{DiccTabla}

% -----------------------------------------------------------------------------
\subsection{Esquema de Interacciones}
\label{subsec:dic_interacciones}

\paragraph{Tabla \texttt{chats} (conversaciones).}
\begin{DiccTabla}{chats}
chat\_id & uuid & NO & \texttt{gen\_random\_uuid()} & PK. \\ \hline
company\_profile\_id & uuid & NO & — & FK a \texttt{profiles\_company} (CASCADE). \\ \hline
basic\_profile\_id & uuid & NO & — & FK a \texttt{profiles\_basic} (CASCADE). \\ \hline
created\_at & timestamptz & NO & \texttt{now()} & Creación. \\ \hline
updated\_at & timestamptz & NO & \texttt{now()} & Última actividad (trigger). \\ \hline
deleted\_at & timestamptz & SÍ & — & Borrado lógico global. \\ \hline
deleted\_at\_basic & timestamptz & SÍ & — & Ocultado por el profesional. \\ \hline
deleted\_at\_company & timestamptz & SÍ & — & Ocultado por el reclutador. \\ \hline
hidden\_basic & boolean & NO & \texttt{false} & Oculto para el profesional. \\ \hline
hidden\_company & boolean & NO & \texttt{false} & Oculto para el reclutador. \\
\end{DiccTabla}

\paragraph{Tabla \texttt{chat\_messages} (mensajes).}
\begin{DiccTabla}{chat\_messages}
message\_id & uuid & NO & \texttt{gen\_random\_uuid()} & PK. \\ \hline
chat\_id & uuid & NO & — & FK a \texttt{chats} (CASCADE). \\ \hline
sender\_id & uuid & NO & — & FK a \texttt{auth.users(id)} (CASCADE). \\ \hline
content & text & NO & \texttt{''} & Contenido del mensaje. \\ \hline
attachment\_type & varchar & SÍ & — & image/audio/file/video (CHECK). \\ \hline
attachment\_url & text & SÍ & — & URL del adjunto. \\ \hline
is\_read & boolean & NO & \texttt{false} & Mensaje leído. \\ \hline
created\_at & timestamptz & NO & \texttt{now()} & Envío. \\ \hline
deleted\_at & timestamptz & SÍ & — & Borrado lógico. \\ \hline
client\_message\_id & uuid & SÍ & — & ID de cliente para idempotencia (único parcial). \\
\end{DiccTabla}

\clearpage
\paragraph{Tabla \texttt{chat\_typing} (indicador de escritura).}
\begin{DiccTabla}{chat\_typing}
chat\_id & uuid & NO & — & PK compuesta. FK a \texttt{chats} (CASCADE). \\ \hline
user\_id & uuid & NO & — & PK compuesta. Usuario que escribe. \\ \hline
updated\_at & timestamptz & NO & \texttt{now()} & Última señal de escritura. \\
\end{DiccTabla}

\paragraph{Tabla \texttt{notifications} (notificaciones).}
\begin{DiccTabla}{notifications}
notification\_id & bigint & NO & (secuencia) & PK. \\ \hline
profile\_id & uuid & NO & — & Destinatario. \\ \hline
type & text & NO & \texttt{'system'} & endorsement/ visit/ recommendation/system /message (CHECK). \\ \hline
title & text & NO & — & Título de la notificación. \\ \hline
message & text & NO & \texttt{''} & Cuerpo. \\ \hline
is\_read & boolean & NO & \texttt{false} & Leída. \\ \hline
created\_at & timestamptz & NO & \texttt{now()} & Creación. \\
\end{DiccTabla}

\paragraph{Tabla \texttt{profile\_likes} (pipeline de candidatos).}
\begin{DiccTabla}{profile\_likes}
like\_id & uuid & NO & \texttt{gen\_random\_uuid()} & PK. \\ \hline
company\_profile\_id & uuid & NO & — & FK a \texttt{profiles\_company} (CASCADE). \\ \hline
basic\_profile\_id & uuid & NO & — & FK a \texttt{profiles\_basic} (CASCADE). \\ \hline
created\_at & timestamptz & NO & \texttt{now()} & Creación. \\ \hline
stage & varchar(20) & NO & \texttt{'saved'} & saved/review/interview/ offer/hired (CHECK). \\ \hline
stage\_order & integer & NO & \texttt{0} & Orden dentro de la fase Kanban. \\ \hline
updated\_at & timestamptz & NO & \texttt{now()} & Modificación. \\
\end{DiccTabla}

\paragraph{Tabla \texttt{recruiter\_candidate\_notes} (notas privadas).}
\begin{DiccTabla}{recruiter\_candidate\_notes}
note\_id & uuid & NO & \texttt{gen\_random\_uuid()} & PK. \\ \hline
company\_profile\_id & uuid & NO & — & FK a \texttt{profiles\_company} (CASCADE). \\ \hline
basic\_profile\_id & uuid & NO & — & FK a \texttt{profiles\_basic} (CASCADE). \\ \hline
body & text & NO & — & Contenido (1–4000 car., CHECK). \\ \hline
created\_at & timestamptz & NO & \texttt{now()} & Creación. \\ \hline
updated\_at & timestamptz & NO & \texttt{now()} & Modificación. \\
\end{DiccTabla}

\paragraph{Tabla \texttt{recruiter\_talent\_views} (vistas de talento).}
\begin{DiccTabla}{recruiter\_talent\_views}
view\_id & uuid & NO & \texttt{gen\_random\_uuid()} & PK. \\ \hline
company\_profile\_id & uuid & NO & — & FK a \texttt{profiles\_company} (CASCADE). \\ \hline
basic\_profile\_id & uuid & NO & — & FK a \texttt{profiles\_basic} (CASCADE). \\ \hline
viewed\_at & timestamptz & NO & \texttt{now()} & Momento de la visita. \\
\end{DiccTabla}

\clearpage
\paragraph{Tabla \texttt{talent\_score\_snapshots} (leaderboard).}
\begin{DiccTabla}{talent\_score\_snapshots}
snapshot\_id & uuid & NO & \texttt{gen\_random\_uuid()} & PK. \\ \hline
basic\_profile\_id & uuid & NO & — & FK a \texttt{profiles\_basic} (CASCADE). \\ \hline
niche & varchar(20) & NO & — & Nicho de talento. \\ \hline
score & numeric(6,2) & NO & \texttt{0} & Puntaje calculado. \\ \hline
rank\_position & integer & SÍ & — & Posición en el ranking. \\ \hline
week\_start & date & NO & — & Semana del snapshot (único por nicho). \\ \hline
created\_at & timestamptz & NO & \texttt{now()} & Creación. \\
\end{DiccTabla}

% -----------------------------------------------------------------------------
\subsection{Esquema de Seguridad y Soporte}
\label{subsec:dic_seguridad}

\paragraph{Tabla \texttt{login\_attempts} (anti fuerza bruta).}
\begin{DiccTabla}{login\_attempts}
id & uuid & NO & \texttt{gen\_random\_uuid()} & PK. \\ \hline
email & text & NO & — & Correo (único por \texttt{lower(email)}). \\ \hline
failed\_count & integer & NO & \texttt{0} & Intentos fallidos consecutivos. \\ \hline
blocked\_until & timestamptz & SÍ & — & Bloqueo temporal (5 fallos → 15 min). \\ \hline
last\_attempt & timestamptz & NO & \texttt{now()} & Último intento. \\ \hline
created\_at & timestamptz & NO & \texttt{now()} & Creación. \\
\end{DiccTabla}

\clearpage
\paragraph{Tabla \texttt{profile\_security\_codes} (OTP de operaciones sensibles).}
\begin{DiccTabla}{profile\_security\_codes}
id & uuid & NO & \texttt{gen\_random\_uuid()} & PK. \\ \hline
profile\_id & uuid & NO & — & FK a \texttt{auth.users(id)} (CASCADE). \\ \hline
code & char(6) & NO & — & Código OTP de 6 dígitos. \\ \hline
purpose & varchar(30) & NO & — & change\_password /change\_ email /delete\_ account /reset\_password (CHECK). \\ \hline
used & boolean & NO & \texttt{false} & Código consumido. \\ \hline
expires\_at & timestamptz & NO & \texttt{now()+15min} & Expiración. \\ \hline
created\_at & timestamptz & NO & \texttt{now()} & Creación. \\ \hline
attempt\_count & integer & NO & \texttt{0} & Intentos fallidos (5 → invalidado). \\
\end{DiccTabla}

\paragraph{Tabla \texttt{verification\_codes} (verificación de cambios).}
\begin{DiccTabla}{verification\_codes}
id & uuid & NO & \texttt{gen\_random\_uuid()} & PK. \\ \hline
auth\_id & uuid & NO & — & FK a \texttt{auth.users(id)} (CASCADE). \\ \hline
code & varchar(6) & NO & — & Código de 6 dígitos (regex \texttt{\^{}\textbackslash d\{6\}\$}, CHECK). \\ \hline
purpose & varchar(30) & NO & — & change\_password /change\_email /delete\_account (CHECK). \\ \hline
new\_value & text & SÍ & — & Valor nuevo solicitado. \\ \hline
expires\_at & timestamptz & NO & — & Expiración (> created\_at, CHECK). \\ \hline
used\_at & timestamptz & SÍ & — & Momento de uso. \\ \hline
created\_at & timestamptz & NO & \texttt{now()} & Creación. \\
\end{DiccTabla}

\clearpage
\paragraph{Tabla \texttt{email\_change\_requests} (cambio de correo).}
\begin{DiccTabla}{email\_change\_requests}
id & uuid & NO & \texttt{gen\_random\_uuid()} & PK. \\ \hline
auth\_id & uuid & NO & — & Usuario solicitante. \\ \hline
new\_email & text & NO & — & Nuevo correo solicitado. \\ \hline
token\_hash & text & NO & — & Hash del token (único). \\ \hline
consumed & boolean & NO & \texttt{false} & Solicitud consumida. \\ \hline
expires\_at & timestamptz & NO & — & Expiración. \\ \hline
created\_at & timestamptz & NO & \texttt{now()} & Creación. \\
\end{DiccTabla}

\begin{AvisoNota}
Las tablas de seguridad (\texttt{login\_attempts}, \texttt{profile\_security\_codes},
\texttt{verification\_codes}, \texttt{email\_change\_requests}) tienen RLS habilitado
pero \textbf{sin políticas para usuarios finales}: solo son accesibles mediante
funciones \texttt{SECURITY DEFINER} y el rol \texttt{service\_role}, blindando los
códigos OTP y los \textit{tokens} frente al acceso directo desde el cliente.
\end{AvisoNota}

% =============================================================================
% 6.3 INTEGRIDAD REFERENCIAL Y RESTRICCIONES
% =============================================================================
\section{Integridad Referencial y Restricciones}
\label{sec:integridad}

La integridad de los datos se garantiza a \textbf{nivel de motor}, no solo en la capa de
aplicación. PostgreSQL valida claves primarias, foráneas, restricciones de unicidad y
de verificación antes de aceptar cualquier escritura, de modo que un cliente
malicioso o un error de la aplicación no puedan corromper el modelo.

% -----------------------------------------------------------------------------
\subsection{Llaves Primarias y Foráneas con sus políticas}
\label{subsec:pk_fk}

\paragraph{Llaves primarias.} El esquema usa mayoritariamente \textbf{UUID v4} como
clave primaria (\texttt{gen\_random\_uuid()}), con dos excepciones notables:
\comando{notifications} usa un \texttt{bigint} secuencial (alto volumen, orden natural),
y \comando{profile\_preferences}, \comando{profiles\_basic}, \comando{profiles\_company}
usan como PK el propio \texttt{id\_auth} (relación 1:1 con \texttt{auth.users}). La tabla
\comando{chat\_typing} emplea una \textbf{PK compuesta} \texttt{(chat\_id, user\_id)}.

\paragraph{Llaves foráneas y política de borrado.} La política de borrado dominante es
\comando{ON DELETE CASCADE}: al eliminar un perfil, se eliminan en cascada todos sus
proyectos, experiencias, habilidades, chats, etc. La única excepción es
\comando{portfolio\_settings.selected\_cv\_document\_id}, que usa \comando{ON DELETE SET
NULL} para no perder la configuración del portafolio si se borra el CV seleccionado. La
figura~\ref{fig:fk-cascada} ilustra el grafo de propagación de borrado desde la raíz.

\clearpage
\vspace{0.3cm}
\begin{figure}[h!]
\centering
\begin{tikzpicture}[
    ent/.style={rectangle, draw=c4Sistema!70!black, fill=c4Componente!25, font=\sffamily\scriptsize,
        align=center, minimum width=2.4cm, minimum height=0.6cm, inner sep=2pt},
    raiz/.style={ent, fill=c4Persona, text=white, font=\sffamily\scriptsize\bfseries},
    casc/.style={-{Stealth[length=2mm]}, semithick, draw=bgFail!80!black},
    setnull/.style={-{Stealth[length=2mm]}, semithick, dashed, draw=colorAcento!80!black}]
    \node[raiz] (auth) at (0,3) {auth.users};
    \node[raiz] (pb) at (-3,1.8) {profiles\_basic};
    \node[raiz] (pc) at (3,1.8) {profiles\_company};

    \node[ent] (c1) at (-5.2,0.4) {projects};
    \node[ent] (c2) at (-3,0.4) {work\_experiences};
    \node[ent] (c3) at (-0.9,0.4) {hard\_skills};
    \node[ent] (c1b) at (-5.2,-0.6) {project\_media};
    \node[ent] (pset) at (-3,-0.6) {portfolio\_settings};
    \node[ent] (cv) at (-0.9,-0.6) {cv\_documents};

    \node[ent] (k1) at (2.4,0.4) {chats};
    \node[ent] (k2) at (4.8,0.4) {profile\_likes};
    \node[ent] (k1b) at (2.4,-0.6) {chat\_messages};

    \draw[casc] (auth) -- node[c4etiqueta]{CASCADE} (pb);
    \draw[casc] (auth) -- node[c4etiqueta]{CASCADE} (pc);
    \foreach \n in {c1,c2,c3} \draw[casc] (pb) -- (\n);
    \draw[casc] (c1) -- (c1b);
    \draw[casc] (pb) -- (pset);
    \draw[casc] (pb) -- (cv);
    \draw[setnull] (cv) to[bend left=20] node[c4etiqueta, right]{SET NULL} (pset);
    \draw[casc] (pc) -- (k1); \draw[casc] (pc) -- (k2);
    \draw[casc] (k1) -- (k1b);
\end{tikzpicture}
\caption{Grafo de propagación de borrado: CASCADE (rojo) y SET NULL (amarillo).}
\label{fig:fk-cascada}
\end{figure}

\begin{NotaTecnica}
Varias FK apuntan directamente a \texttt{auth.users(id)} en lugar de a la tabla de
perfil (\texttt{profile\_connections}, \texttt{profile\_preferences},
\texttt{profile\_security\_codes}, \texttt{verification\_codes}, \texttt{chat\_messages.sender\_id}).
Esto ancla esos datos a la identidad de autenticación, garantizando consistencia con
las políticas RLS que comparan contra \texttt{auth.uid()}.
\end{NotaTecnica}

\paragraph{Restricciones de unicidad.} Más allá de las PK, el esquema define UNIQUE de
negocio resumidas en la tabla~\ref{tab:unique}.

\vspace{0.2cm}
\noindent
\renewcommand{\arraystretch}{1.3}
\fontsize{8.5}{10.2}\selectfont\sffamily
\begin{tabularx}{\textwidth}{|>{\ttfamily\color{colorPrimario}}l|X|}
\hline
\rowcolor{headerTabla}
\sffamily\textbf{Tabla} & \sffamily\textbf{Restricción de unicidad} \\
\hline
chats & \texttt{(company\_profile\_id, basic\_profile\_id)}: un único chat por par. \\
\hline
\rowcolor{grisTecnico}
profile\_likes & \texttt{(company\_profile\_id, basic\_profile\_id)}: un like por par. \\
\hline
recruiter\_talent\_views & \texttt{(company\_profile\_id, basic\_profile\_id)}: una vista por par. \\
\hline
\rowcolor{grisTecnico}
portfolio\_settings & \texttt{(profile\_id)}: un portafolio por perfil. \\
\hline
global\_skill\_tags & \texttt{(name)}: nombre de habilidad único. \\
\hline
\rowcolor{grisTecnico}
talent\_score\_snapshots & \texttt{(basic\_profile\_id, niche, week\_start)}: un snapshot por nicho/semana. \\
\hline
portfolio\_items & \texttt{(portfolio\_settings\_id, item\_type, item\_id)}: ítem único por portafolio. \\
\hline
\rowcolor{grisTecnico}
email\_change\_requests & \texttt{(token\_hash)}: token de cambio único. \\
\hline
\end{tabularx}
\label{tab:unique}

% -----------------------------------------------------------------------------
\subsection{Restricciones de verificación (CHECK)}
\label{subsec:check}

Las restricciones \textbf{CHECK} validan dominios de datos atómicos directamente en el
motor, actuando como una capa de defensa independiente de la validación de la
aplicación. La tabla~\ref{tab:checks} cataloga las más relevantes; nótese cómo
codifican las máquinas de estados y enumeraciones del dominio.

\clearpage
\vspace{0.2cm}
\noindent
\renewcommand{\arraystretch}{1.3}
\fontsize{8.2}{9.8}\selectfont\sffamily
\begin{tabularx}{\textwidth}{|>{\ttfamily\color{colorPrimario}}l|>{\ttfamily}l|X|}
\hline
\rowcolor{headerTabla}
\sffamily\textbf{Tabla} & \sffamily\textbf{Columna} & \sffamily\textbf{Dominio validado (CHECK)} \\
\hline
projects & status & draft, in\_progress, completed, archived \\
\hline
\rowcolor{grisTecnico}
projects & category & Web, Mobile, API, Data, DevOps, Other \\
\hline
projects & visibility & Public, Private \\
\hline
\rowcolor{grisTecnico}
hard\_skills & level & Junior, Mid, Senior \\
\hline
hard\_skills & top\_order & NULL o entero 1–3; requiere \texttt{is\_top = true} \\
\hline
\rowcolor{grisTecnico}
profile\_likes & stage & saved, review, interview, offer, hired \\
\hline
academic\_records & education\_type & university, master\_degree, phd, certification, course, bootcamp, high\_school \\
\hline
\rowcolor{grisTecnico}
academic\_records & gpa & NULL o rango 0–100 \\
\hline
profile\_connections & provider & email, github, google, linkedin, … (+20), custom \\
\hline
\rowcolor{grisTecnico}
profile\_connections & api\_health & healthy, degraded, down \\
\hline
chat\_messages & attachment\_type & image, audio, file, video \\
\hline
\rowcolor{grisTecnico}
notifications & type & endorsement, visit, recommendation, system, message \\
\hline
verification\_codes & code & exactamente 6 dígitos (regex \texttt{\^{}\textbackslash d\{6\}\$}) \\
\hline
\rowcolor{grisTecnico}
work\_experiences & (varios) & fechas coherentes; descripción <= 1000 car.; campos no vacíos \\
\hline
\end{tabularx}
\label{tab:checks}

\begin{AvisoNota}
Las restricciones de fecha del tipo \comando{end\_date >= start\_date} aparecen de forma
consistente en \texttt{projects}, \texttt{work\_experiences} y \texttt{academic\_records},
evitando rangos temporales inválidos a nivel de motor. El CHECK
\comando{chk\_top\_order\_requires\_is\_top} es un ejemplo de regla de negocio compuesta
embebida en la base de datos.
\end{AvisoNota}

% =============================================================================
% 6.4 ESTRATEGIA DE INDEXACIÓN Y OPTIMIZACIÓN
% =============================================================================
\section{Estrategia de Indexación y Optimización de Consultas}
\label{sec:indexacion}

El esquema define alrededor de \textbf{80 índices} (incluyendo los implícitos de PK y
UNIQUE). La estrategia no es indexar indiscriminadamente, sino atacar los patrones de
consulta reales del backend: filtrado por dueño (\texttt{profile\_id}), exclusión de
borrados lógicos y ordenación por relevancia o fecha.

% -----------------------------------------------------------------------------
\subsection{Índices B-Tree sobre identificadores y llaves foráneas}
\label{subsec:btree}

Todas las llaves foráneas de uso frecuente cuentan con un índice B-Tree explícito que
acelera los \textit{joins} y los filtros por dueño. La tabla~\ref{tab:tipos-indice}
clasifica los índices del esquema por tipo y propósito.

\vspace{0.2cm}
\noindent
\renewcommand{\arraystretch}{1.3}
\fontsize{8.5}{10.2}\selectfont\sffamily
\begin{tabularx}{\textwidth}{|>{\bfseries\color{colorPrimario}}l|c|X|}
\hline
\rowcolor{headerTabla}
\textbf{Tipo de índice} & \textbf{Aprox.} & \textbf{Propósito y ejemplo} \\
\hline
B-Tree simple (FK) & $\sim$20 & Acelera \textit{joins} y filtros por dueño. Ej.: \texttt{idx\_projects\_profile\_id}. \\
\hline
\rowcolor{grisTecnico}
Compuesto & $\sim$12 & Cubre filtro + orden en una pasada. Ej.: \texttt{idx\_notifications\_profile (profile\_id, is\_read, created\_at DESC)}. \\
\hline
Parcial (\textit{soft-delete}) & $\sim$10 & Indexa solo filas activas (\texttt{WHERE deleted\_at IS NULL}), reduciendo tamaño. \\
\hline
\rowcolor{grisTecnico}
Único parcial & $\sim$6 & Unicidad condicional. Ej.: \texttt{uq\_profile\_tag} único por \texttt{(profile\_id, tag\_id)} solo si activo. \\
\hline
Funcional & 1 & \texttt{idx\_login\_attempts\_email} sobre \texttt{lower(email)} para búsqueda case-insensitive. \\
\hline
\end{tabularx}
\label{tab:tipos-indice}

\subsection{Índices compuestos y parciales para búsquedas optimizadas}
\label{subsec:compuestos}

Los \textbf{índices parciales} son la pieza central de la estrategia de
\textit{soft-delete}: como la mayoría de las consultas filtran por
\texttt{deleted\_at IS NULL}, indexar solo las filas activas reduce el tamaño del índice
y mejora la selectividad. La tabla~\ref{tab:indices-clave} muestra ejemplos
representativos extraídos del esquema real.

\vspace{0.2cm}
\noindent
\renewcommand{\arraystretch}{1.25}
\fontsize{8}{9.6}\selectfont\sffamily
\begin{tabularx}{\textwidth}{|>{\ttfamily\color{colorPrimario}\scriptsize}l|X|}
\hline
\rowcolor{headerTabla}
\sffamily\textbf{Índice} & \sffamily\textbf{Definición y motivación} \\
\hline
idx\_projects\_profile\_active & \texttt{(profile\_id, is\_featured DESC, created\_at DESC) WHERE deleted\_at IS NULL}. Lista los proyectos activos de un perfil, destacados primero. \\
\hline
\rowcolor{grisTecnico}
idx\_work\_experiences\_active & \texttt{(profile\_id, is\_current DESC, start\_date DESC) WHERE deleted\_at IS NULL}. Línea de tiempo laboral. \\
\hline
idx\_notifications\_profile & \texttt{(profile\_id, is\_read, created\_at DESC)}. Bandeja de notificaciones del usuario. \\
\hline
\rowcolor{grisTecnico}
idx\_profile\_likes\_company\_stage & \texttt{(company\_profile\_id, stage, stage\_order)}. Render del pipeline Kanban por fase. \\
\hline
uq\_chat\_messages\_client\_id & \texttt{(chat\_id, client\_message\_id) WHERE client\_message\_id IS NOT NULL}. Garantiza idempotencia de mensajes. \\
\hline
\rowcolor{grisTecnico}
idx\_portfolio\_settings\_slug & \texttt{(slug) WHERE slug IS NOT NULL AND slug <> ''}. Resolución de portafolios públicos por slug. \\
\hline
idx\_psc\_active\_lookup & \texttt{(profile\_id, purpose, code) WHERE used = false}. Validación rápida de OTP vigentes. \\
\hline
\rowcolor{grisTecnico}
idx\_tss\_niche\_week & \texttt{(niche, week\_start, rank\_position)}. Cálculo del leaderboard por nicho. \\
\hline
\end{tabularx}
\label{tab:indices-clave}

\begin{NotaTecnica}
El índice \comando{uq\_chat\_messages\_client\_id} merece especial atención: al ser un
\textbf{índice único parcial} sobre \texttt{client\_message\_id}, implementa la
idempotencia de mensajes a nivel de motor. Si el cliente reintenta un envío con el mismo
\texttt{client\_message\_id}, el \texttt{INSERT} viola la unicidad y se evita el
duplicado, sin necesidad de lógica adicional en la aplicación.
\end{NotaTecnica}

% =============================================================================
% 6.5 LÓGICA PROGRAMABLE EN EL MOTOR DE BASE DE DATOS
% =============================================================================
\section{Lógica Programable en el Motor de Base de Datos}
\label{sec:programable}

EthosHub traslada una parte significativa de la lógica de negocio sensible al propio
motor PostgreSQL, mediante \textbf{\textit{triggers}} (disparadores) y
\textbf{funciones PL/pgSQL}. Esta decisión persigue dos objetivos: garantizar
invariantes que no deben depender de la aplicación (auditoría, exclusividad) y operar de
forma segura por encima de las políticas RLS mediante funciones \texttt{SECURITY DEFINER}.

% -----------------------------------------------------------------------------
\subsection{Triggers para auditoría y sincronización de eventos}
\label{subsec:triggers}

El esquema define 12 \textit{triggers}, clasificables en cuatro familias funcionales,
resumidas en la tabla~\ref{tab:triggers}.

\clearpage
\vspace{0.2cm}
\noindent
\renewcommand{\arraystretch}{1.25}
\fontsize{8.2}{9.8}\selectfont\sffamily
\begin{tabularx}{\textwidth}{|>{\ttfamily\color{colorPrimario}\scriptsize}l|>{\scriptsize}l|X|}
\hline
\rowcolor{headerTabla}
\sffamily\textbf{Trigger} & \sffamily\textbf{Tabla / Evento} & \sffamily\textbf{Función} \\
\hline
notify\_chat\_message & chat\_messages / AFTER INSERT & Crea una notificación para el destinatario del mensaje. \\
\hline
\rowcolor{grisTecnico}
trg\_chat\_message\_timestamp & chat\_messages / AFTER INSERT & Actualiza \texttt{updated\_at} del chat (orden de conversaciones). \\
\hline
notify\_profile\_like & profile\_likes / AFTER INSERT & Notifica al candidato que ha recibido un \textit{like}. \\
\hline
\rowcolor{grisTecnico}
trg\_cv\_documents\_updated\_at & cv\_documents / BEFORE UPDATE & Mantiene \texttt{updated\_at} en cada modificación. \\
\hline
trg\_profile\_connections\_updated\_at & profile\_connections / BEFORE UPDATE & Actualiza \texttt{updated\_at}. \\
\hline
\rowcolor{grisTecnico}
trg\_check\_basic\_exclusivity & profiles\_basic / BEFORE INSERT & Impide que un usuario sea profesional y reclutador a la vez. \\
\hline
trg\_check\_company\_exclusivity & profiles\_company / BEFORE INSERT & Misma exclusividad para el perfil reclutador. \\
\hline
\rowcolor{grisTecnico}
log\_profile\_basic\_events & profiles\_basic / INSERT,UPDATE & Audita altas y borrados en \texttt{admin.system\_logs}. \\
\hline
log\_project\_created & projects / AFTER INSERT & Registro de auditoría de creación de proyectos. \\
\hline
\rowcolor{grisTecnico}
log\_portfolio\_visibility & portfolio\_settings / UPDATE & Audita cambios de publicación del portafolio. \\
\hline
log\_login\_blocked & login\_attempts / UPDATE & Audita bloqueos por fuerza bruta. \\
\hline
\rowcolor{grisTecnico}
log\_skill\_tag\_created & global\_skill\_tags / AFTER INSERT & Audita la creación de nuevas etiquetas de habilidad. \\
\hline
\end{tabularx}
\label{tab:triggers}

El siguiente \textit{trigger} ejemplifica la \textbf{sincronización de eventos}: al
insertarse un mensaje, genera automáticamente una notificación para el destinatario,
resolviendo quién es el receptor a partir de la conversación.

\begin{lstlisting}[language=SQL, caption={Trigger de notificación de mensaje (\texttt{core.notify\_chat\_message}).}]
CREATE FUNCTION core.notify_chat_message() RETURNS trigger
  LANGUAGE plpgsql SECURITY DEFINER SET search_path TO 'core','pg_temp' AS $$
DECLARE v_recipient UUID;
BEGIN
  SELECT CASE WHEN c.basic_profile_id = NEW.sender_id THEN c.company_profile_id
              ELSE c.basic_profile_id END
  INTO v_recipient FROM core.chats c WHERE c.chat_id = NEW.chat_id;

  IF v_recipient IS NOT NULL AND v_recipient <> NEW.sender_id THEN
    INSERT INTO core.notifications (profile_id, type, title, message)
    VALUES (v_recipient, 'message', 'Nuevo mensaje',
            left(COALESCE(NULLIF(NEW.content,''), 'Has recibido un adjunto'), 120));
  END IF;
  RETURN NEW;
END; $$;
\end{lstlisting}

\paragraph{Auditoría histórica SCD Tipo 2.} El esquema \comando{history} implementa
\textbf{Slowly Changing Dimensions de Tipo 2} sobre los perfiles: cada cambio relevante
genera una nueva versión histórica (con marcas de vigencia), preservando el rastro
completo de la evolución de cada perfil en \texttt{history.hist\_profiles} y
\texttt{history.hist\_profiles\_basic}.

% -----------------------------------------------------------------------------
\subsection{Procedimientos almacenados y funciones PL/pgSQL}
\label{subsec:funciones}

El esquema expone alrededor de \textbf{60 funciones}, la mayoría declaradas como
\comando{SECURITY DEFINER} para operar de forma controlada por encima de RLS. Se
agrupan funcionalmente como muestra la tabla~\ref{tab:rpcs}.

\vspace{0.2cm}
\noindent
\renewcommand{\arraystretch}{1.25}
\fontsize{8.2}{9.8}\selectfont\sffamily
\begin{tabularx}{\textwidth}{|>{\bfseries\color{colorPrimario}}l|X|}
\hline
\rowcolor{headerTabla}
\textbf{Grupo} & \textbf{Funciones representativas (RPC)} \\
\hline
Chat & \texttt{send\_message}, \texttt{get\_messages}, \texttt{mark\_messages\_read}, \texttt{upsert\_chat}, \texttt{delete\_chat}, \texttt{get\_chat\_list}, \texttt{get\_unread\_count}. \\
\hline
\rowcolor{grisTecnico}
Pipeline / Likes & \texttt{add\_like}, \texttt{remove\_like}, \texttt{update\_like\_stage}, \texttt{get\_company\_likes}, \texttt{has\_liked\_profile}. \\
\hline
Talento & \texttt{record\_talent\_view}, \texttt{get\_recent\_talent\_views}, \texttt{get\_explore\_portfolios}, \texttt{get\_profile\_likes\_count}. \\
\hline
\rowcolor{grisTecnico}
Seguridad & \texttt{check\_login\_block}, \texttt{record\_login\_failure}, \texttt{clear\_login\_attempts}, \texttt{increment\_otp\_attempt}, \texttt{cleanup\_expired\_codes}. \\
\hline
Identidad & \texttt{handle\_new\_user}, \texttt{delete\_profile\_cascade}, \texttt{sync\_auth\_email\_to\_profiles}, \texttt{sync\_oauth\_identity\_to\_connections}. \\
\hline
\rowcolor{grisTecnico}
Administración & \texttt{admin.search\_profiles}, \texttt{admin.get\_global\_metrics}, \texttt{admin.soft\_delete\_profile}, \texttt{admin.restore\_profile}, \texttt{admin.log\_event}. \\
\hline
\end{tabularx}
\label{tab:rpcs}

La función \comando{record\_login\_failure} ilustra la lógica anti fuerza bruta
ejecutada de forma atómica en el motor: incrementa el contador, reinicia si el bloqueo
expiró y aplica un bloqueo de 15 minutos al quinto fallo, todo en una sola sentencia
\texttt{UPSERT}.

\begin{lstlisting}[language=SQL, caption={Función anti fuerza bruta (\texttt{core.record\_login\_failure}, fragmento).}]
INSERT INTO core.login_attempts (email, failed_count, last_attempt)
VALUES (lower(p_email), 1, now())
ON CONFLICT (lower(email)) DO UPDATE
  SET failed_count = CASE
        WHEN login_attempts.blocked_until < now() THEN 1          -- bloqueo expirado: reinicia
        ELSE login_attempts.failed_count + 1 END,
      blocked_until = CASE
        WHEN (...failed_count + 1) >= 5                            -- 5 intentos
        THEN now() + '15 minutes'::INTERVAL ELSE NULL END,
      last_attempt = now()
RETURNING *;
\end{lstlisting}

\begin{AlertaSeguridad}
El uso de \comando{SECURITY DEFINER} es potente pero delicado: estas funciones se
ejecutan con los privilegios de su propietario, eludiendo RLS. Por ello \textbf{fijan
explícitamente} \comando{SET search\_path} (p.~ej. \texttt{'core','pg\_temp'}) para
prevenir ataques de secuestro de \textit{search path}, y validan internamente la
pertenencia del llamante antes de operar sobre datos ajenos.
\end{AlertaSeguridad}

% =============================================================================
% 6.6 GESTIÓN DE TRANSACCIONES Y CONTROL DE CONCURRENCIA
% =============================================================================
\section{Gestión de Transacciones y Control de Concurrencia}
\label{sec:concurrencia}

PostgreSQL emplea \textbf{MVCC} (\textit{Multi-Version Concurrency Control}): los lectores
nunca bloquean a los escritores y viceversa, ya que cada transacción ve una instantánea
consistente de los datos. EthosHub se apoya en este modelo y lo complementa con
estrategias específicas en la capa de persistencia.

\subsection{Niveles de aislamiento transaccional}
\label{subsec:aislamiento}

El backend opera bajo el nivel de aislamiento por defecto de PostgreSQL,
\comando{READ COMMITTED}, adecuado para la mayoría de las operaciones CRUD. La
tabla~\ref{tab:aislamiento} resume las garantías relevantes.

\vspace{0.2cm}
\noindent
\renewcommand{\arraystretch}{1.3}
\fontsize{9}{10.8}\selectfont\sffamily
\begin{tabularx}{\textwidth}{|>{\bfseries\color{colorPrimario}}l|X|}
\hline
\rowcolor{headerTabla}
\textbf{Aspecto} & \textbf{Comportamiento en EthosHub} \\
\hline
Nivel por defecto & \texttt{READ COMMITTED}: cada sentencia ve los datos confirmados al iniciarse. \\
\hline
\rowcolor{grisTecnico}
Idempotencia & El índice único parcial de \texttt{chat\_messages} evita duplicados ante reintentos concurrentes. \\
\hline
Operaciones atómicas & \texttt{UPSERT} (\texttt{INSERT ... ON CONFLICT}) en \texttt{add\_like}, \texttt{record\_login\_failure}, \texttt{upsert\_chat}. \\
\hline
\rowcolor{grisTecnico}
Pooling & El acceso pasa por el \textit{pooler} de Supabase, acotando conexiones simultáneas (RNF-CONC-01). \\
\hline
\end{tabularx}
\label{tab:aislamiento}

\subsection{Control transaccional desde la capa de persistencia}
\label{subsec:control_tx}

Desde el backend, la capa tipada con \textbf{jOOQ} y \textbf{Spring JDBC} delimita las
transacciones; las operaciones compuestas (p.~ej. crear un perfil y su configuración)
se ejecutan dentro de una única transacción gestionada por Spring
(\texttt{@Transactional}). Las funciones \texttt{SECURITY DEFINER} que realizan varias
escrituras (como \comando{delete\_profile\_cascade}) constituyen, por sí mismas, una
unidad transaccional atómica en el motor.

\begin{NotaTecnica}
El patrón \comando{INSERT ... ON CONFLICT DO NOTHING/UPDATE} se prefiere a la secuencia
``leer-decidir-escribir'' porque resuelve la condición de carrera \textbf{dentro de una
única sentencia atómica}, eliminando ventanas de tiempo en las que dos transacciones
concurrentes podrían insertar duplicados.
\end{NotaTecnica}

% =============================================================================
% 6.7 SEGURIDAD DEL MOTOR Y POLÍTICAS RLS
% =============================================================================
\section{Seguridad del Motor y Políticas RLS}
\label{sec:rls}

La seguridad de los datos no descansa únicamente en el backend: PostgreSQL aplica
\textbf{Row Level Security} (RLS) directamente en el motor, de modo que incluso un acceso
directo desde el cliente web (con la \texttt{anon key} de Supabase) queda restringido a
las filas que el usuario tiene derecho a ver. Esta es una defensa en profundidad: aunque
fallara una comprobación en la aplicación, el motor seguiría protegiendo los datos.

\subsection{Configuración de Políticas de Seguridad a Nivel de Fila (RLS)}
\label{subsec:politicas_rls}

De las 25 tablas de \texttt{core}, \textbf{23 tienen RLS habilitado}. Las dos excepciones
(\comando{portfolio\_items} y \\ \comando{talent\_score\_snapshots}) se acceden
exclusivamente a través del backend o de funciones, nunca directamente desde el cliente.
La figura~\ref{fig:rls-flujo} ilustra los dos caminos de acceso a los datos y cómo RLS
gobierna el acceso directo del cliente.

\clearpage
\vspace{0.3cm}
\begin{figure}[h!]
\centering
\begin{tikzpicture}[node distance=0.9cm and 1.6cm,
    box/.style={rectangle, rounded corners=3pt, draw, font=\sffamily\scriptsize, align=center,
        minimum width=2.8cm, minimum height=0.9cm, inner sep=3pt},
    rel/.style={-{Stealth[length=2mm]}, semithick}]
    \node[box, fill=c4Componente!40] (cli) at (0,2.5) {Cliente SPA\\\scriptsize (anon key)};
    \node[box, fill=c4Componente!40] (be) at (0,0) {Backend\\\scriptsize (service\_role / JWT)};
    \node[box, fill=colorAcento!30, draw=colorAcento!80!black] (rls) at (4.4,2.5) {Políticas RLS\\\scriptsize auth.uid() = dueño};
    \node[box, fill=c4Datos!25, draw=c4Datos!70!black] (db) at (4.4,0) {PostgreSQL\\\scriptsize esquema core};
    \node[box, fill=c4Componente!40] (rpc) at (4.4,-1.6) {RPC SECURITY\\DEFINER};

    \draw[rel] (cli) -- node[c4etiqueta, above]{SELECT directo} (rls);
    \draw[rel] (rls) -- node[c4etiqueta, right]{filas permitidas} (db);
    \draw[rel] (be) -- node[c4etiqueta, left]{jOOQ} (db);
    \draw[rel] (be) to[bend right=20] node[c4etiqueta, below]{invoca} (rpc);
    \draw[rel] (rpc) -- (db);
\end{tikzpicture}
\caption{Doble vía de acceso: el cliente pasa por RLS; el backend opera por jOOQ y RPCs.}
\label{fig:rls-flujo}
\end{figure}

\paragraph{Patrones de política.} Las $\sim$70 políticas RLS se reducen a unos pocos
patrones reutilizados, resumidos en la tabla~\ref{tab:patrones-rls}.

\vspace{0.2cm}
\noindent
\renewcommand{\arraystretch}{1.3}
\fontsize{8.2}{9.8}\selectfont\sffamily
\begin{tabularx}{\textwidth}{|>{\bfseries\color{colorPrimario}}l|>{\ttfamily\scriptsize}l|X|}
\hline
\rowcolor{headerTabla}
\textbf{Patrón} & \textbf{Expresión} & \textbf{Aplicación} \\
\hline
Dueño directo & profile\_id = auth.uid() & academic\_records, work\_experiences, cv\_documents, profile\_connections. \\
\hline
\rowcolor{grisTecnico}
Participante de chat & EXISTS(... chats ...) & chat\_messages: solo los dos participantes ven/insertan mensajes. \\
\hline
Lectura pública & visibility = 'Public' & projects, project\_media: lectura por \texttt{anon} solo de contenido público no borrado. \\
\hline
\rowcolor{grisTecnico}
Portafolio publicado & is\_published = true & portfolio\_settings: visible si está publicado. \\
\hline
Doble parte & company OR basic = uid() & profile\_likes, chats: visible para reclutador o candidato. \\
\hline
\rowcolor{grisTecnico}
Solo service\_role & true (service\_role) & verification\_codes, hard\_skills, soft\_skills: gestionados por el backend. \\
\hline
\end{tabularx}
\label{tab:patrones-rls}

A continuación, dos políticas reales que ejemplifican los patrones más sofisticados: la
visibilidad de mensajes restringida a los participantes del chat, y la nueva política
(v1.1.0) que permite a un candidato ver los \textit{likes} recibidos.

\begin{lstlisting}[language=SQL, caption={Políticas RLS representativas (esquema \texttt{core}).}]
-- Solo los participantes del chat ven sus mensajes
CREATE POLICY "participants can view messages" ON core.chat_messages FOR SELECT
USING (EXISTS (SELECT 1 FROM core.chats c
   WHERE c.chat_id = chat_messages.chat_id
     AND (c.company_profile_id = auth.uid() OR c.basic_profile_id = auth.uid())));

-- El candidato puede ver los likes que ha recibido (migración V15)
CREATE POLICY profile_likes_owner_select ON core.profile_likes FOR SELECT
USING (basic_profile_id = auth.uid());
\end{lstlisting}

\subsection{Gestión de roles de base de datos}
\label{subsec:roles}

El modelo de acceso de Supabase define tres roles principales sobre los que se evalúan
las políticas RLS. La tabla~\ref{tab:roles} describe cada uno y su uso en EthosHub.

\clearpage
\vspace{0.2cm}
\noindent
\renewcommand{\arraystretch}{1.3}
\fontsize{9}{10.8}\selectfont\sffamily
\begin{tabularx}{\textwidth}{|>{\ttfamily\bfseries\color{colorPrimario}}l|X|}
\hline
\rowcolor{headerTabla}
\sffamily\textbf{Rol} & \sffamily\textbf{Descripción y uso} \\
\hline
anon & Usuario \textbf{no autenticado}. Solo accede a lectura pública (portafolios publicados, proyectos \texttt{Public}). \\
\hline
\rowcolor{grisTecnico}
authenticated & Usuario \textbf{autenticado} (JWT válido). Las políticas comparan \texttt{auth.uid()} contra el dueño de la fila. \\
\hline
service\_role & Rol \textbf{privilegiado} del backend. Evade RLS (\texttt{USING true}) para operaciones de servidor; su clave nunca se expone al cliente. \\
\hline
\end{tabularx}
\label{tab:roles}

\begin{AlertaSeguridad}
La \comando{anon key} es \textbf{pública por diseño} y se expone en el cliente; su
seguridad reposa íntegramente en RLS. La \comando{service\_role key}, en cambio, evade
RLS y \textbf{nunca} debe llegar al navegador: vive solo en el backend como variable de
entorno. La separación estricta entre ambas claves es la piedra angular del modelo de
seguridad de datos de EthosHub.
\end{AlertaSeguridad}

\subsection{Síntesis de la cobertura de seguridad}
\label{subsec:sintesis_rls}

La tabla~\ref{tab:cobertura-rls} ofrece una visión consolidada de la cobertura de RLS
por tabla, evidenciando que ninguna tabla con datos de usuario queda desprotegida.

\vspace{0.2cm}
\noindent
\renewcommand{\arraystretch}{1.2}
\fontsize{8}{9.6}\selectfont\sffamily
\begin{tabularx}{\textwidth}{|>{\ttfamily\scriptsize\color{colorPrimario}}l|c|c|X|}
\hline
\rowcolor{headerTabla}
\sffamily\textbf{Tabla} & \sffamily\textbf{RLS} & \sffamily\textbf{Pol.} & \sffamily\textbf{Modelo de acceso} \\
\hline
profiles\_basic & Sí & 2 & Dueño (escritura) + lectura autenticada. \\ \hline
profiles\_company & Sí & 2 & Dueño + lectura autenticada. \\ \hline
\rowcolor{grisTecnico}
projects & Sí & 3 & Dueño + lectura pública + service\_role. \\ \hline
project\_media & Sí & 3 & Dueño vía proyecto + público + service\_role. \\ \hline
\rowcolor{grisTecnico}
work\_experiences & Sí & 4 & Dueño (CRUD completo). \\ \hline
academic\_records & Sí & 4 & Dueño (CRUD completo). \\ \hline
\rowcolor{grisTecnico}
hard\_skills / soft\_skills & Sí & 2 & Lectura autenticada + service\_role. \\ \hline
cv\_documents & Sí & 1 & Dueño (ALL). \\ \hline
\rowcolor{grisTecnico}
portfolio\_settings & Sí & 2 & Dueño + publicado. \\ \hline
profile\_connections & Sí & 2 & Dueño + service\_role. \\ \hline
\rowcolor{grisTecnico}
chats & Sí & 7 & Participantes (con borrado por parte). \\ \hline
chat\_messages & Sí & 7 & Participantes (ver/insertar/borrar propios). \\ \hline
\rowcolor{grisTecnico}
profile\_likes & Sí & 8 & Reclutador (CRUD) + candidato (ver recibidos). \\ \hline
recruiter\_candidate\_notes & Sí & 1 & Reclutador dueño. \\ \hline
\rowcolor{grisTecnico}
recruiter\_talent\_views & Sí & 4 & Reclutador dueño. \\ \hline
notifications & Sí & 2 & Dueño (ver/actualizar). \\ \hline
\rowcolor{grisTecnico}
verification\_codes & Sí & 1 & Solo service\_role (OTP blindado). \\ \hline
portfolio\_items & No & 0 & Acceso solo vía backend. \\ \hline
\rowcolor{grisTecnico}
talent\_score\_snapshots & No & 0 & Acceso solo vía funciones. \\ \hline
\end{tabularx}
\label{tab:cobertura-rls}

\clearpage
\begin{AvisoNota}
La tabla \comando{profile\_likes} concentra 8 políticas porque modela un recurso
compartido entre dos actores con derechos asimétricos: el reclutador gestiona el
\textit{like} (crear, mover de fase, eliminar), mientras que el candidato solo puede
\textbf{ver} los que ha recibido. Este matiz, introducido en la migración V15, es un
ejemplo de cómo RLS expresa reglas de negocio finas directamente en el motor.
\end{AvisoNota}

